\chapter{Introduction}

\section{Background}

Malaria remains one of the most pressing public health challenges in sub-Saharan Africa. According to the World Health Organization \citep{who2023malaria}, there were an estimated 249 million malaria cases and 608,000 malaria deaths worldwide in 2022. The African region bears the heaviest burden, accounting for approximately 94\% of all malaria cases and 95\% of all malaria deaths. To put this in perspective, one child dies from malaria approximately every minute in Africa.

Children under five years of age are particularly vulnerable, comprising nearly 80\% of all malaria deaths in the region. This vulnerability arises from several factors: immature immune systems that have not yet developed immunity to malaria parasites, higher rates of severe anemia when infected, and limited access to healthcare in many rural areas. The economic burden is also substantial, with malaria-related absenteeism and treatment costs pushing millions of families deeper into poverty.

Kenya has a heterogeneous malaria transmission landscape. The western region near Lake Victoria experiences high perennial transmission, with entomological inoculation rates (EIR) reaching 100-300 infective bites per person per year. In contrast, the central highlands around Nairobi and the northern arid areas have low or seasonal transmission, with EIR below 10 infective bites per person per year \citep{noor2009risks}. This heterogeneity is driven by altitude, rainfall patterns, temperature, and land use. Despite significant progress, malaria remains a leading cause of morbidity and mortality among Kenyan children under five, accounting for approximately 20\% of all outpatient visits in high-transmission areas.

Insecticide-treated nets (ITNs) are among the most widely implemented malaria prevention interventions. ITNs create a physical barrier between humans and mosquitoes while also killing or repelling mosquitoes through insecticide (typically pyrethroids). The insecticide component is critical because it reduces mosquito longevity, killing mosquitoes before they can transmit parasites, and provides a "mass effect" when coverage is high. Large-scale distribution campaigns have increased ITN ownership across Africa, with the proportion of households owning at least one ITN increasing from 5\% in 2004 to 56\% in 2021 \citep{who2023malaria}. However, ownership does not guarantee use; studies consistently show usage rates 10-20 percentage points lower than ownership rates.

\section{The Empirical Challenge}

While ITNs have been shown effective in randomized controlled trials \citep{lengeler2004insecticide, pryce2018insecticide}, estimating their effect using observational data presents significant challenges due to confounding. Randomization ensures that, in expectation, treated and control groups are balanced on all characteristics, both observed and unobserved. Observational studies lack this protection.

Key sources of confounding include:

\begin{enumerate}
	\item \textbf{Socioeconomic status}: Wealthier households are more likely to own and use ITNs. They also tend to live in better-quality housing with fewer mosquito entry points (screened windows, closed eaves, concrete walls). They have better nutrition, which improves immune response to malaria. They have better access to healthcare, reducing malaria morbidity and mortality. Wealth is therefore a confounder that, if uncontrolled, biases estimates away from the true effect.
	
	\item \textbf{Environmental conditions}: Rainfall creates mosquito breeding sites. Temperature affects the parasite development rate within mosquitoes (the extrinsic incubation period). Vegetation influences mosquito habitat and resting sites. These factors affect both malaria transmission intensity and ITN use decisions (people may use nets more when they perceive higher mosquito density).
	
	\item \textbf{Demographic factors}: Child age affects both ITN use patterns (caregivers are more likely to ensure ITN use for younger children) and malaria immunity (older children have had more exposure, potentially developing partial immunity). Sex may affect ITN use through behavioral differences.
	
	\item \textbf{Household factors}: Electricity access is a proxy for wealth and also affects ITN use through access to television malaria prevention messages. Water and sanitation affect mosquito breeding sites near the home.
	
	\item \textbf{Health-seeking behavior}: Families that prioritize malaria prevention may also seek prompt treatment when fever occurs, reducing severe outcomes.
\end{enumerate}

To obtain unbiased estimates, these confounders must be adequately controlled. This study uses a combination of multivariable adjustment, propensity score methods, and Double Machine Learning to address confounding.

\section{Research Gap}

Despite extensive literature on ITNs and malaria, several important gaps remain.

\textbf{First}, most existing studies report associations rather than effects estimated under explicit causal assumptions. The difference is not merely semantic. An associational study might report that "ITN users have 40\% lower malaria risk than non-users." A causal study makes explicit the assumptions under which this difference can be interpreted as the effect of ITN use. These assumptions include conditional ignorability (no unmeasured confounding), positivity (every child has some chance of using an ITN), and no interference (one child's ITN use does not affect another child's risk). This study explicitly states and, where possible, tests these assumptions.

\textbf{Second}, many analyses ignore the hierarchical structure of survey data where children are nested within clusters (enumeration areas). Preliminary analysis reveals an intra-cluster correlation (ICC) of 50-58\%, indicating that more than half of the variation in malaria risk is between clusters. Ignoring this clustering leads to underestimated standard errors (overconfident inferences) and inefficient estimates. A child in western Kenya has fundamentally different baseline malaria risk than a child in the Nairobi highlands, regardless of ITN use. This study accounts for clustering using GLMM with random intercepts and cluster-aware cross-fitting in DML.

\textbf{Third}, existing research typically estimates a single type of effect without distinguishing between cluster-specific (conditional) and population-average (marginal) effects. As \citet{neuhaus1991comparison} demonstrate, when ICC is high, conditional effects are typically larger in magnitude than marginal effects. The conditional effect answers: "For a child in a typical cluster, how much does ITN use reduce their risk?" The marginal effect answers: "If we gave every child in Kenya an ITN, what would be the average reduction in malaria risk across the country?" These are different questions requiring different methods and leading to different answers. Policymakers need both: the conditional effect to understand biological efficacy, and the marginal effect to predict population impact.

\section{Research Question}

This study addresses the following research question:

\begin{pro}
	What is the effect of insecticide-treated net (ITN) use on malaria infection among children under five in Kenya?
\end{pro}

We estimate two types of effects. First, \textbf{cluster-specific (conditional) effects}: the effect within a typical cluster, holding cluster-level heterogeneity constant. Second, \textbf{population-average (marginal) effects}: the average effect across the entire population.

Secondary questions include:
\begin{enumerate}
	\item How do cluster-specific estimates compare with population-average estimates?
	\item Does the impact vary across wealth quintiles, residence type, age groups, and rainfall levels?
	\item How robust are the estimates to different machine learning learners and cross-fitting folds?
\end{enumerate}

\section{Motivation for Causal Machine Learning}

Traditional causal inference methods rely on parametric models that assume linearity in the log-odds and may miss important non-linear relationships. For example, the relationship between age and malaria risk is not linear: infants have some protection from maternal antibodies, risk peaks around 12-23 months as maternal antibodies wane, then gradually declines as acquired immunity develops. Linear models cannot capture this pattern.

Double Machine Learning \citep{chernozhukov2018double} combines machine learning algorithms with causal inference techniques, offering several advantages:

\begin{enumerate}
	\item \textbf{Flexibility}: Machine learning learners can capture non-linearities and interactions without requiring the researcher to specify the functional form a priori.
	
	\item \textbf{High-dimensional confounding}: DML handles many covariates with regularization, potentially including interactions and non-linear terms.
	
	\item \textbf{Double robustness}: The DML estimator remains consistent if either the outcome model or the treatment model is correctly specified.
	
	\item \textbf{Valid inference}: Cross-fitting procedures ensure that standard errors are valid despite using flexible machine learning methods.
	
	\item \textbf{Neyman orthogonality}: The score function is constructed so that small errors in estimating nuisance functions do not create first-order bias.
\end{enumerate}

We use DML2 (Definition 3.2 in \citealp{chernozhukov2018double}) with 5-fold and 10-fold cross-fitting.

\section{Study Contributions}

\subsection{Methodological Contributions}

This study applies several methods in a novel combination for the Kenyan context:

\begin{enumerate}
	\item \textbf{Explicit distinction between conditional and marginal effects}: We estimate both using appropriate methods (GLMM and PSM for conditional, DML and DR for marginal) and compare them.
	
	\item \textbf{GLMM for propensity scores}: Most studies estimate propensity scores using logistic regression, ignoring clustering. We extend this to GLMM following \citet{ayele2024assessment}, who demonstrated the importance of multilevel modeling for malaria outcomes in Kenya.
	
	\item \textbf{Cluster-aware cross-fitting for DML}: Standard cross-fitting can split children from the same cluster across different folds, creating information leakage. We implement cluster-aware folds where entire clusters stay together.
	
	\item \textbf{Triangulation of six methods}: We compare results across logistic regression, GLMM, PSM, IPTW, DR, and DML to assess convergence.
	
	\item \textbf{Integration of high-resolution environmental data}: We link survey clusters to DHS Geocov environmental data (rainfall, temperature, EVI, aridity, wet days) using GPS coordinates.
\end{enumerate}

\subsection{Policy Contributions}

The study provides rigorous evidence on ITN effectiveness through:

\begin{enumerate}
	\item \textbf{Causal estimates}: Unlike associational studies, we provide estimates that can be interpreted as effects under explicit assumptions.
	
	\item \textbf{Heterogeneity analysis}: We identify subgroups where ITNs are most effective (poorer wealth quintile, older children aged 36-59 months, high-rainfall areas, urban areas), informing targeted distribution strategies.
	
	\item \textbf{Documentation of declining ITN use}: ITN use decreased from 58.0\% in 2015 to 50.5\% in 2020. This decline occurred despite continued distribution campaigns, suggesting that distribution alone is insufficient and behavior change communication is needed.
\end{enumerate}

\section{Thesis Structure}

This thesis is organized as follows. Chapter 2 reviews the existing literature on ITN effectiveness, causal inference methods, and plagiarism. Chapter 3 presents the methodology, including data sources, variables, causal framework, and complete mathematical derivations of all statistical methods. Chapter 4 presents the empirical results with all tables and figures. Chapter 5 discusses the findings, interprets each figure, and addresses limitations. Chapter 6 concludes with a summary of key findings, answers to the research question, limitations, contributions, and future research directions.

\section{Scope and Delimitations}

This study focuses on children under five in Kenya using the 2015 and 2020 Kenya Malaria Indicator Surveys (KMIS). Several delimitations apply:

\begin{enumerate}
	\item \textbf{Excluded surveys}: The 2014 and 2022 Demographic and Health Surveys (DHS) were excluded because they lack the malaria outcome variable (rapid diagnostic test results).
	
	\item \textbf{Complete case analysis}: The study uses complete case analysis, with missingness below 5\% for all covariates, following \citet{little2019statistical}. This assumes missingness is completely at random (MCAR).
	
	\item \textbf{Causal assumptions}: Interpretation relies on conditional ignorability and positivity. Following \citet{westreich2010invited}, positivity requires that for every combination of observed confounders, there is a non-zero probability of both ITN use and non-use. We assess this empirically in Chapter 4.
	
	\item \textbf{Generalizability}: Findings are specific to Kenya and may not generalize to other countries with different transmission intensities, mosquito species, or ITN distribution strategies.
	
	\item \textbf{Unmeasured confounding}: Unmeasured confounders may still bias estimates. This limitation is addressed using E-values sensitivity analysis in Chapter 4 and discussed in Chapter 5.
\end{enumerate}